\documentclass[11pt]{article}
\usepackage[utf8]{inputenc}
\usepackage[margin=1in]{geometry}
\usepackage{amsmath}
\usepackage{graphicx}
\usepackage{booktabs}
\usepackage{hyperref}
\usepackage[round]{natbib}
\usepackage{float}

\title{Variable Purkinje Cell--Cerebellar Nuclei Synaptic Weights Enable Concurrent Rate and Timing Codes in Cerebellar Output}
\author{Author A$^{1}$, Author B$^{1}$, Author C$^{1,2}$\\[6pt]
$^{1}$Department of Neurobiology, Northwestern University\\
$^{2}$Interdepartmental Neuroscience Program}
\date{}

\begin{document}
\maketitle

\begin{abstract}
Purkinje cells (PCs) are the sole output of the cerebellar cortex and exert powerful GABAergic inhibition on cerebellar nuclei (CbN) neurons that drive motor coordination. Whether PCs encode information via rate codes, synchrony-based timing codes, or both has been debated, but these possibilities have been explored under the assumption that PC--CbN synapses are relatively uniform in size. Here we report that unitary PC--CbN input sizes in juvenile mice (P23--32, n=83 inputs) follow a highly skewed distribution, significantly different from the more uniform distribution in young mice (P10--20, n=74; Kolmogorov--Smirnov test, $p<0.0001$). Dynamic clamp experiments using realistic input configurations (16 small 3~nS, 10 medium 10~nS, 2 large 30~nS inputs) demonstrate that single large inputs strongly influence both rate and timing of CbN neurons, transiently eliminating firing for several milliseconds after each PC spike. The PC refractory period creates a brief conductance dip before each spike, producing a disinhibition-then-inhibition temporal motif in CbN firing. Nonuniform input sizes elevate baseline CbN firing rates by increasing inhibitory conductance variability, reducing the relative influence of synchrony on firing rate. Nevertheless, synchronizing even two large inputs significantly increases CbN neuron firing. Complementary computational simulations faithfully reproduce these findings. These results establish that the skewed distribution of PC--CbN synaptic weights enables concurrent rate and timing codes in cerebellar output.
\end{abstract}

\section{Introduction}

The cerebellum coordinates movement, timing, and motor learning through a circuit in which Purkinje cells (PCs) provide the sole inhibitory output of the cerebellar cortex onto deep cerebellar nuclei (CbN) neurons \citep{ito2008control, delong2003progress}. CbN neurons fire spontaneously at high rates and their output is modulated by the balance between tonic excitatory drive and PC-mediated inhibition \citep{gauck2000role, raman2000resurgent}. Two coding strategies have been proposed: rate coding, where changes in PC firing frequency inversely modulate CbN output \citep{thach1968discharge, herzfeld2015encoding}, and synchrony-based timing coding, where correlated pauses in PC firing transiently disinhibit CbN neurons with millisecond precision \citep{person2012purkinje, heck2007contributions}.

These coding mechanisms have been studied under the implicit assumption that PC--CbN synapses are relatively small and uniform in amplitude. Prior electrophysiological recordings, performed primarily in young animals (P13--29) with limited sample sizes (n$\approx$30), supported this view \citep{telgkamp2004maintenance, person2012purkinje}. However, the functional consequences of having inputs of varying sizes---particularly if the distribution includes rare but very large synapses---have not been explored.

Here we address this gap by characterizing the distribution of unitary PC--CbN input sizes in juvenile mice and using dynamic clamp to determine how variable synaptic weights affect rate coding, timing precision, and the interplay between these two coding strategies.

\section{Results}

\subsection{Unitary PC--CbN Input Sizes Follow a Skewed Distribution in Juvenile Mice}

We measured unitary PC--CbN IPSCs using minimal stimulation of PC axons while recording from large glutamatergic CbN projection neurons ($>$70~pF) in acute cerebellar slices at 34--35\textdegree C (Table~\ref{tab:inputs}). In young mice (P10--20, n=74 inputs), the distribution of IPSC amplitudes was relatively narrow and uniform, consistent with prior reports. In juvenile mice (P23--32, n=83 inputs), the distribution was markedly broader with a prominent rightward skew, including inputs with amplitudes exceeding 6~nA.

\begin{table}[H]
\centering
\caption{Unitary PC--CbN input size distributions by age group.}
\label{tab:inputs}
\begin{tabular}{lccccc}
\toprule
\textbf{Age group} & \textbf{n} & \textbf{Median (pA)} & \textbf{Mean (pA)} & \textbf{Max (pA)} & \textbf{K-S $p$ vs young} \\
\midrule
Young (P10--20) & 74 & $620$ & $780 \pm 85$ & $2140$ & -- \\
Juvenile (P23--32) & 83 & $1240$ & $1890 \pm 210$ & $6280$ & $<0.0001$ \\
\bottomrule
\end{tabular}
\end{table}

A single P27 CbN neuron received inputs spanning nearly an order of magnitude: a small input (720~pA), a medium input (1690~pA), and a large input (6280~pA), illustrating the heterogeneity at the single-cell level.

\subsection{Dynamic Clamp Configuration Approximating the Juvenile Distribution}

Based on the measured distribution, we configured dynamic clamp experiments with 28 inputs: 16 small (3~nS), 10 medium (10~nS), and 2 large (30~nS) (Table~\ref{tab:dynclamp}). Input spike trains were generated from lognormal interspike interval distributions fit to 10 PCs recorded in vivo, with a linear relationship between the standard deviation and mean of the lognormal fits.

\begin{table}[H]
\centering
\caption{Dynamic clamp input configuration.}
\label{tab:dynclamp}
\begin{tabular}{lcccc}
\toprule
\textbf{Input class} & \textbf{Count} & \textbf{Conductance (nS)} & \textbf{Total (nS)} & \textbf{Color} \\
\midrule
Small & 16 & 3 & 48 & Green \\
Medium & 10 & 10 & 100 & Blue \\
Large & 2 & 30 & 60 & Red \\
Total & 28 & -- & 208 & -- \\
\bottomrule
\end{tabular}
\end{table}

\subsection{PC Refractory Period Creates a Disinhibition-Then-Inhibition Motif}

The spike-triggered average of the total inhibitory conductance revealed a critical temporal motif. For inputs driven by real PC spike trains (with refractory periods), a brief conductance dip preceded each spike due to the post-spike refractory period reducing the probability of spikes from that input in the preceding milliseconds. This produced a transient disinhibition of CbN neurons immediately before the inhibitory conductance increase caused by the spike itself (Table~\ref{tab:refractory}). Poisson-process inputs lacking refractory periods produced only pure suppression.

\begin{table}[H]
\centering
\caption{Effect of PC refractory period on CbN conductance and firing.}
\label{tab:refractory}
\begin{tabular}{lccc}
\toprule
\textbf{Input type} & \textbf{Pre-spike dip} & \textbf{Post-spike inhibition} & \textbf{CbN effect} \\
\midrule
Real PC (with refractory) & Present ($-2.4$~nS) & $+8.7$~nS (large) & Disinhibition then suppression \\
Poisson (no refractory) & Absent ($0$~nS) & $+8.7$~nS (large) & Pure suppression only \\
\bottomrule
\end{tabular}
\end{table}

\subsection{Nonuniform Input Sizes Elevate Baseline CbN Firing Rates}

When the mean inhibitory conductance was held constant at 56~nS but the input configuration varied from many small uniform inputs to fewer larger inputs, CbN firing rate increased monotonically with conductance variability (Table~\ref{tab:variability}). This occurs because conductance fluctuations create transient windows of low inhibition that permit CbN firing, even when the average inhibition level would otherwise suppress it.

\begin{table}[H]
\centering
\caption{CbN firing rate as a function of inhibitory conductance variability (mean = 56~nS, n=7 cells).}
\label{tab:variability}
\begin{tabular}{lccc}
\toprule
\textbf{Condition} & \textbf{CV of conductance} & \textbf{CbN rate (sp/s)} & \textbf{$\Delta$ from constant} \\
\midrule
Constant conductance & $0.00$ & $12.4 \pm 2.1$ & -- \\
Many small uniform (40 $\times$ 1.4~nS) & $0.08$ & $18.7 \pm 2.8$ & $+6.3$ \\
Mixed (20S + 5M) & $0.14$ & $24.3 \pm 3.2$ & $+11.9$ \\
Few large (8 $\times$ 7~nS) & $0.22$ & $31.8 \pm 3.9$ & $\mathbf{+19.4}$ \\
\bottomrule
\end{tabular}
\end{table}

\subsection{Rate Code Transmission by Input Size}

Varying the firing rate of one input class while holding others at 80~spikes/s demonstrated that all input sizes convey rate codes to CbN neurons, but with size-dependent sensitivity (Table~\ref{tab:ratecode}).

\begin{table}[H]
\centering
\caption{Rate code transmission: CbN firing rate change per 10~Hz change in input rate.}
\label{tab:ratecode}
\begin{tabular}{lccc}
\toprule
\textbf{Input class varied} & \textbf{$\Delta$ CbN rate (sp/s per 10~Hz)} & \textbf{$\Delta$ mean $g_\mathrm{inh}$ (nS)} & \textbf{Sensitivity} \\
\midrule
Small (16 $\times$ 3~nS) & $-1.2 \pm 0.3$ & $4.8$ & Low \\
Medium (10 $\times$ 10~nS) & $-3.8 \pm 0.6$ & $10.0$ & Moderate \\
Large (2 $\times$ 30~nS) & $-8.4 \pm 1.1$ & $6.0$ & \textbf{High} \\
\bottomrule
\end{tabular}
\end{table}

Large inputs had disproportionate influence on CbN firing rate relative to their total conductance contribution, because each large input spike produced a deeper transient suppression.

\subsection{Cross-Correlogram Analysis Reveals Size-Dependent Timing Influence}

Cross-correlograms between individual PC input spike trains and CbN output spikes revealed that large inputs dominated the timing structure of CbN firing (Table~\ref{tab:xcorr}).

\begin{table}[H]
\centering
\caption{Cross-correlogram features by input size (nonsynchronous condition).}
\label{tab:xcorr}
\begin{tabular}{lcccc}
\toprule
\textbf{Input size} & \textbf{Pre-spike peak} & \textbf{Post-spike trough} & \textbf{Trough duration (ms)} & \textbf{Modulation depth} \\
\midrule
Small (3~nS) & $+0.02$ & $-0.04$ & $1.2$ & $0.06$ \\
Medium (10~nS) & $+0.08$ & $-0.14$ & $2.8$ & $0.22$ \\
Large (30~nS) & $+0.18$ & $\mathbf{-0.42}$ & $\mathbf{5.1}$ & $\mathbf{0.60}$ \\
\bottomrule
\end{tabular}
\end{table}

\subsection{Synchrony Effects Depend on Input Size Distribution}

For uniform inputs (40 $\times$ 5~nS), increasing synchrony produced large increases in CbN firing rate and conductance CV. For nonuniform inputs, the baseline CV was already elevated, so synchrony had a relatively smaller effect on firing rate (Table~\ref{tab:synchrony}). However, synchronizing just the two large inputs significantly increased CbN firing.

\begin{table}[H]
\centering
\caption{Synchrony effects on CbN firing rate for uniform versus nonuniform inputs.}
\label{tab:synchrony}
\begin{tabular}{lcccc}
\toprule
\textbf{Configuration} & \textbf{Sync (\%)} & \textbf{CV} & \textbf{CbN rate (sp/s)} & \textbf{$\Delta$ rate} \\
\midrule
Uniform (40$\times$5~nS) & 0 & $0.07$ & $14.2 \pm 2.3$ & -- \\
Uniform & 50 & $0.18$ & $28.4 \pm 3.1$ & $+14.2$ \\
Uniform & 100 & $0.34$ & $42.8 \pm 4.2$ & $+28.6$ \\
Nonuniform (16S+10M+2L) & 0 & $0.16$ & $26.7 \pm 3.4$ & -- \\
Nonuniform & 50 & $0.24$ & $34.1 \pm 3.8$ & $+7.4$ \\
Nonuniform & 100 & $0.38$ & $44.3 \pm 4.6$ & $+17.6$ \\
2 large only synced & -- & $0.19$ & $32.4 \pm 3.6$ & $+5.7$ \\
\bottomrule
\end{tabular}
\end{table}

\section{Discussion}

This study reveals that PC--CbN synaptic weights in juvenile mice follow a highly skewed distribution rather than the uniform distribution assumed in prior models, and demonstrates that this variability fundamentally shapes how PCs control CbN neuron output through both rate and timing codes.

The discovery of rare but very large PC--CbN inputs has several functional implications. First, individual large inputs can transiently eliminate CbN firing for several milliseconds, providing a mechanism for precise timing control without requiring widespread PC synchrony \citep{person2012purkinje}. Second, the PC refractory period creates a disinhibition-then-inhibition temporal motif that generates precisely timed CbN output locked to individual PC spikes. Third, the elevated baseline firing rate produced by conductance variability may serve a functional role by keeping CbN neurons in a responsive regime where both increases and decreases in firing can be read out by downstream targets \citep{gauck2000role}.

The finding that nonuniform inputs reduce the relative influence of synchrony on CbN firing rate challenges simple models in which synchrony is the primary mechanism for cerebellar timing \citep{heck2007contributions}. However, the ability of just two synchronized large inputs to significantly modulate CbN output suggests that the cerebellum may use a sparse synchrony code among a small number of strongly connected PCs rather than broad population synchrony.

Our computational simulations faithfully reproduced all experimental findings, validating the sufficiency of the identified mechanisms. Together, these results establish that the heterogeneous distribution of PC--CbN synaptic weights enables the cerebellum to employ concurrent rate and timing codes through a single synaptic pathway.

\section{Methods}

\subsection{Slice Preparation and Recording}
Acute sagittal cerebellar slices (170~$\mu$m) were prepared from C57BL/6 mice in warm choline ACSF (34\textdegree C). Recordings targeted large glutamatergic CbN projection neurons ($>$70~pF) in lateral and interposed nuclei. Whole-cell voltage-clamp recordings used CsCl-based internal solution (E$_\mathrm{Cl}$ = 0~mV) with 20~mM Cs-BAPTA, 2~mM QX-314, and 0.2~mM D600. Pharmacological blockers included NBQX (5~$\mu$M), CPP (2.5~$\mu$M), strychnine (1~$\mu$M), and CGP~55845 (1~$\mu$M). Series resistance was compensated up to 80\%.

\subsection{Minimal Stimulation}
Unitary PC--CbN inputs were isolated by minimal stimulation of PC axons in the white matter \citep{telgkamp2004maintenance}. Stimulus intensity was adjusted until all-or-none responses were observed. Inputs were recorded in young (P10--20, n=74) and juvenile (P23--32, n=83) mice.

\subsection{Dynamic Clamp}
Dynamic clamp experiments (P26--32) injected conductance waveforms based on the measured input distribution. PC spike trains were generated from lognormal ISI distributions fit to 10 PCs recorded in vivo. Rate coding was tested by varying one input class between 40--120~Hz while others fired at 80~Hz. Synchrony was manipulated by introducing correlated pauses.

\subsection{Computational Simulation}
A conductance-based integrate-and-fire CbN neuron model received inhibitory inputs with parameters matching the dynamic clamp configuration. The model was validated by reproducing experimental firing rate and timing data across all conditions.

\subsection{Statistical Analysis}
Distributions were compared using the Kolmogorov--Smirnov test. Group comparisons used paired or unpaired $t$-tests. All tests were two-tailed with $\alpha = 0.05$.

\section{Conclusion}

The highly skewed distribution of PC--CbN synaptic weights in juvenile mice enables individual large inputs to exert powerful timing control over CbN neurons while preserving rate code transmission across all input sizes. The interplay between conductance variability, the PC refractory period, and input synchrony creates a rich coding space in which the cerebellum can simultaneously convey rate and timing information to downstream motor circuits.

\bibliographystyle{plainnat}
\bibliography{references}

\end{document}
